\documentclass[a4paper,12pt]{article}
\usepackage[T2A]{fontenc}
\usepackage[utf8]{inputenc}
\usepackage[russian]{babel}
\usepackage{amsmath,amssymb,mathtools}
\usepackage{geometry}
\geometry{left=22mm,right=22mm,top=22mm,bottom=23mm}
\usepackage{helvet}
\renewcommand{\familydefault}{\sfdefault}
\usepackage{microtype}
\usepackage{parskip}
\usepackage{booktabs,tabularx,array}
\usepackage{xcolor}
\usepackage{tikz}
\usepackage{hyperref}
\usepackage{fancyhdr}
\definecolor{accent}{HTML}{007F84}
\definecolor{deep}{HTML}{23464B}
\definecolor{muted}{HTML}{46545A}
\hypersetup{colorlinks=true,linkcolor=accent,urlcolor=accent,pdfauthor={Лёвин Артём Александрович},pdftitle={Ревью домашней работы. Ксения Клыкова. 27 октября 2026}}
\setlength{\parindent}{0pt}
\setlength{\parskip}{7pt}
\linespread{1.13}
\setlength{\headheight}{14pt}
\pagestyle{fancy}
\fancyhf{}
\fancyhead[L]{\small\color{muted}Математика / Делимость и делители}
\fancyhead[R]{\small\color{muted}Ксения Клыкова}
\fancyfoot[C]{\small\thepage}
\renewcommand{\headrulewidth}{0.25pt}
\newcommand{\keyidea}{\textcolor{accent}{\textbf{Ключевая идея.}}\ }
\newcommand{\errorhead}{\textcolor{red!65!black}{\textbf{Ошибка.}}\ }
\newcommand{\result}{\textcolor{deep}{\textbf{Итог.}}\ }
\newcommand{\tasktitle}[2]{%
  \clearpage\phantomsection\label{task:#1}%
  {\Large\bfseries\color{deep}Задание №\,#1}\par\vspace{2mm}
  {\color{muted}\large #2}\par\vspace{5mm}
}
\begin{document}
\thispagestyle{empty}
\begin{center}
\vspace*{36mm}
{\large\color{accent}МАТЕМАТИКА}\par
\vspace{19mm}
{\LARGE\color{deep}Ревью домашней работы}\par
\vspace{16mm}
{\Large\color{accent}Количество делителей\\[4pt]и каноническое разложение}\par
\vspace{18mm}
{\normalsize Разбор записей и уточнение математических обоснований}\par
\vfill
{\large Ксения Клыкова}\par
\vspace{10mm}
{27 октября 2026 г.}\par
\vspace{15mm}
{\small Лёвин Артём Александрович\\[3pt]эксклюзивно для Ксении Клыковой}\par
\vspace{18mm}
\begin{tikzpicture}
\draw[accent,line width=.85pt] (-4,0) .. controls (-2.3,.65) and (-1.25,-.65) .. (0,0)
 .. controls (1.3,.62) and (2.55,-.64) .. (4,0);
\end{tikzpicture}
\end{center}
\clearpage
{\Large\bfseries\color{deep}Сводная таблица}\par
\vspace{4mm}
\small
\renewcommand{\arraystretch}{1.3}
\begin{tabularx}{\linewidth}{@{}>{\raggedright\arraybackslash}p{13mm}>{\raggedright\arraybackslash}p{31mm}X>{\raggedright\arraybackslash}p{30mm}>{\raggedright\arraybackslash}p{31mm}@{}}
\toprule
\textbf{№} & \textbf{Тема} & \textbf{Суть ошибки} & \textbf{Ваш ответ} & \textbf{Правильный ответ}\\
\midrule
\hyperref[task:8]{8} & Делители числа 72 & В явном списке записано 27 вместо 72; перепутана запись числа и количества делителей. & 12; в списке 27 & 12; в списке 72\\
\addlinespace[5pt]
\hyperref[task:10]{10} & Формула $\tau(N)$ & В записи множителей появились лишние верхние индексы и обозначение $N$ вместо $\tau(N)$. & $\alpha=4$ & $\alpha=4$; корректная формула $\tau(N)=6(\alpha+1)$\\
\bottomrule
\end{tabularx}
\normalsize
\vspace{7mm}
\textcolor{muted}{В таблицу включены только задания с неточностями в математической записи или проверке. Номера ведут к подробному разбору.}
\clearpage
\thispagestyle{empty}
\vspace*{.38\textheight}
\begin{center}
{\Huge\bfseries\color{deep}Разбор ошибок}\par
\vspace{9mm}
{\large\color{muted}Точные места исправлений и методы самопроверки}
\end{center}

\tasktitle{8}{Количество и полный список делителей числа 72}
\textbf{Условие.} Найдите количество натуральных делителей числа $72$ двумя способами: с помощью формулы $\tau$ и явным перечислением.

\textbf{Ваше решение.} Получено число $12$. При этом записано равенство вида $72=(3+1)(2+1)=12$, а при перечислении делителей последним указан $27$ вместо $72$.

\textbf{Ваш ответ.} $12$ делителей.

\errorhead В явном списке присутствует число $27$, которое не делит $72$, и отсутствует сам делитель $72$.

\textbf{Где именно возникает ошибка.} При подсчёте слагаемых-показателей получен верный результат $12$. Первое содержательное расхождение возникает в списке: после $36$ записано $27$. Кроме того, количество делителей следует обозначать $\tau(72)$, поскольку само число $72$ не равно $12$.

\textbf{Почему это неверно.} Натуральное число $d$ --- делитель числа $72$ тогда и только тогда, когда $72:d$ --- целое число. Но
\[
 \frac{72}{27}=\frac83
\]
целым числом не является. Напротив, $72:72=1$, поэтому $72$ обязательно входит в список своих натуральных делителей. Совпадение правильного количества с количеством перечисленных чисел ещё не доказывает, что все перечисленные числа являются делителями.

\keyidea Разложение $72$ на простые множители даёт одновременно формулу количества делителей и способ получить их без пропусков.

\textbf{Исправление от правильного шага.} Используйте верное разложение
\[
72=2^3\cdot3^2.
\]
Значит,
\[
\tau(72)=(3+1)(2+1)=4\cdot3=12.
\]
Из списка необходимо убрать $27$ и вписать $72$.

\newpage\textbf{Полное правильное решение.} Первый способ --- формула. В натуральном делителе числа $72$ показатель двойки можно выбрать четырьмя способами: $0,1,2,3$; показатель тройки --- тремя способами: $0,1,2$. Поэтому всего $4\cdot3=12$ разных делителей.

Второй способ --- явное построение чисел $2^i3^j$:
\[
\begin{array}{c|ccc}
 &j=0&j=1&j=2\\\hline
 i=0&1&3&9\\
 i=1&2&6&18\\
 i=2&4&12&36\\
 i=3&8&24&72
\end{array}
\]
Расположим их в порядке возрастания:
\[
\boxed{1,\ 2,\ 3,\ 4,\ 6,\ 8,\ 9,\ 12,\ 18,\ 24,\ 36,\ 72.}
\]
\textbf{Проверка.} В таблице четыре строки по три числа, то есть $12$ делителей. Каждое из них имеет вид $2^i3^j$ с $0\le i\le3$ и $0\le j\le2$, поэтому действительно делит $72$. Других натуральных делителей у числа $72$ быть не может.

\result $\tau(72)=12$. Последний делитель в возрастающем списке --- $72$.

\textbf{Задача на закрепление.} Найдите количество натуральных делителей числа $108$ с помощью формулы $\tau$, затем выпишите все делители в порядке возрастания.

\tasktitle{10}{Число делителей произведения простых степеней}
\textbf{Условие.} Известно, что $N=2^{\alpha}\cdot3^2\cdot5$, причём $\tau(N)=30$. Найдите натуральное число $\alpha$.

\textbf{Ваше решение.} Из разложения верно получено выражение $6(\alpha+1)$, а затем уравнение
\[
6\alpha+6=30,\qquad 6\alpha=24,\qquad\alpha=4.
\]
В промежуточной записи вместо $\tau(N)$ стоит $N$, а над числами $2$ и $1$ в множителях формулы присутствуют дополнительные верхние индексы.

\textbf{Ваш ответ.} $\alpha=4$.

\errorhead Некорректно оформлен переход от разложения числа на простые множители к формуле для \emph{количества} делителей.

\textbf{Где именно возникает ошибка.} Исходное разложение $N=2^{\alpha}\cdot3^2\cdot5$ и конечное уравнение $6(\alpha+1)=30$ правильны. Проблема возникает в строке, где произведение факторов $\alpha+1$, $2+1$, $1+1$ подписано как $N$ и снабжено лишними надстрочными символами. Если считать их степенями, выражение уже не равно $6(\alpha+1)$.

\textbf{Почему это неверно.} Формула
\[
\tau(p^aq^br^c)=(a+1)(b+1)(c+1)
\]
определяет \emph{количество делителей}, а не само число. Здесь $p,q,r$ --- разные простые числа, $a,b,c$ --- их показатели. В формуле каждое слагаемое $a+1$, $b+1$, $c+1$ берётся в скобки; дополнительно возводить показатели в степени не требуется. Например, $2^3+1=9$, тогда как нужный множитель $2+1=3$.

\newpage\keyidea В формулу $\tau$ подставляют только \emph{показатели} степеней из канонического разложения и к каждому прибавляют единицу.

\textbf{Исправление от правильного шага.} Запишите $5=5^1$ и отдельно вычислите количество делителей:
\[
\tau(N)=(\alpha+1)(2+1)(1+1)=6(\alpha+1).
\]
\textbf{Полное правильное решение.} По условию
\[
N=2^\alpha\cdot3^2\cdot5^1.
\]
У каждого натурального делителя $N$ показатель двойки может принимать значения от $0$ до $\alpha$ (всего $\alpha+1$ вариантов), показатель тройки --- $0,1,2$ (три варианта), показатель пятёрки --- $0,1$ (два варианта). Следовательно,
\[
\tau(N)=(\alpha+1)\cdot3\cdot2=6(\alpha+1).
\]
По условию делителей ровно тридцать:
\begin{align*}
6(\alpha+1)&=30,\\
\alpha+1&=5,\\
\boxed{\alpha}&=4.
\end{align*}
\textbf{Проверка.} При $\alpha=4$ получаем $N=2^4\cdot3^2\cdot5$ и
\[
\tau(N)=(4+1)(2+1)(1+1)=5\cdot3\cdot2=30.
\]
Количество делителей совпадает с условием.

\result $\alpha=4$. Вычисление верное; важна точная запись $\tau(N)$ и показателей в формуле.

\textbf{Задача на закрепление.} Пусть $M=2^{\beta}\cdot3^3\cdot7^2$ и $\tau(M)=48$. Найдите натуральное число $\beta$.
\end{document}